\begin{pastexam}[label={px:ncku-112-prog-09}]{成大資訊聯招 112 學年度 第 9 題}
（選擇）

{[5\%]} Suppose we have a Patricia trie as shown below. Let us insert 0110 into and then delete 1100 from
the trie. What does node 1101 point to?

\begin{center}
\begin{tikzpicture}[x=1cm,y=1cm,>={Stealth[length=1.8mm]},
  k/.style={ellipse,draw,inner sep=0pt,minimum width=11.5mm,minimum height=7mm,font=\sffamily\small},
  w/.style={font=\sffamily\footnotesize,inner sep=1pt}]
\node[k] (h) at (0,0) {0001};
\node[k] (a) at (-1.3,-1.12) {1101};
\node[k] (b) at (-3.0,-2.4) {0011};
\node[k] (c) at (0.1,-2.57) {1001};
\node[k] (d) at (-1.47,-4.28) {1000};
\node[k] (e) at (1.22,-4.44) {1100};
\node[w] at (0.72,0.5) {0};
\draw[->] (h)--node[w,below right,pos=0.75]{1}(a);
\draw[->] (a)--node[w,below right,pos=0.78]{3}(b);
\draw[->] (a)--(c);
\node[w] at (0.83,-2.06) {2};
\draw[->] (c)--node[w,below right,pos=0.6]{4}(d);
\draw[->] (c)--(e);
\node[w] at (1.97,-3.92) {4};

\draw[->] (b.south) .. controls +(0.2,-0.35) and +(0.1,-0.5) .. (-1.9,-2.7) .. controls +(-0.05,0.25) and +(0.4,0) .. (b.east);
\draw[->] (b.south west) .. controls +(-0.7,-0.1) and +(-0.1,-1.2) .. (-3.95,-1.3) .. controls +(0.1,0.9) and +(-1.6,0) .. (h.west);

\draw[->] (d.south west) .. controls +(-0.5,-0.2) and +(0,-0.35) .. (-2.6,-4.55) .. controls +(0,0.3) and +(-0.4,0) .. (d.west);
\draw[->] (d.east) .. controls +(0.7,0) and +(0,-0.9) .. (c.south);

\draw[->] (e.south west) .. controls +(-0.5,-0.2) and +(0,-0.35) .. (0.1,-4.7) .. controls +(0,0.3) and +(-0.4,0) .. (e.west);
\draw[->] (e.east) .. controls +(0.15,1.9) and +(1.9,0) .. (a.east);
\end{tikzpicture}
\end{center}
(A)~0001 \qquad (B)~0110 \qquad (C)~0011 \qquad (D)~1001 \qquad (E)~1000
\end{pastexam}

\begin{pastsoln}
答案：(B)(D)。 $1101$ 的左指標指向 $0110$，右指標指向 $1001$。

\textbf{讀圖}。 照資結書第~11.5.3~節的 Patricia（Horowitz 的版本）：位元從左邊第 $1$ 位數起，標頭 $0001$ 的
$\mathit{bit} = 0$、只用左指標；邊上的數字是指到的那個節點的 $\mathit{bit}$。 第 $\mathit{bit}$ 位是 $0$ 往左、
$1$ 往右；指到 $\mathit{bit}$ 不比自己大的節點的是往回的指標。 由圖讀出每個節點的兩個指標，並用「往回的指標
一定指向第 $\mathit{bit}$ 位相符的鍵」核對左右：
\begin{center}
\small
\begin{tabular}{cccc}
\toprule
節點（鍵） & $\mathit{bit}$ & 左指標 & 右指標 \\
\midrule
$0001$（標頭） & $0$ & $1101$ & — \\
$1101$ & $1$ & $0011$ & $1001$ \\
$0011$ & $3$ & $0001$（往回） & $0011$（自己） \\
$1001$ & $2$ & $1000$ & $1100$ \\
$1000$ & $4$ & $1000$（自己） & $1001$（往回） \\
$1100$ & $4$ & $1100$（自己） & $1101$（往回） \\
\bottomrule
\end{tabular}
\end{center}
例如 $0011$ 看第 $3$ 位：$0011$ 自己第 $3$ 位是 $1$，所以指向自己的是右指標；$0001$ 第 $3$ 位是 $0$，在左邊。

\textbf{插入 $0110$}（資結書虛擬碼~11.12）。
\begin{enumerate}
\item 搜尋：標頭 $\to 1101$（$\mathit{bit}$ 1，$0110$ 第 $1$ 位 $0$，往左）$\to 0011$（$\mathit{bit}$ 3，第 $3$ 位
$1$，往右）$\to$ 指回 $0011$ 自己，停。 $y = 0011$。
\item $0110$ 和 $0011$：第 $1$ 位都是 $0$，第 $2$ 位 $1 \ne 0$，$j = 2$。
\item 從標頭重新往下：$1101$ 的 $\mathit{bit} = 1 < 2$，繼續往左；左指標指向 $0011$，$\mathit{bit} = 3 \ge 2$，停。
新節點 $z$（鍵 $0110$，$\mathit{bit} = 2$）插在 $1101$ 的左指標上。 $0110$ 第 $2$ 位是 $1$，所以 $z$ 的右指標
指向自己，左指標接原本的 $0011$ 子樹。 $1101$ 的左指標改成 $z$。
\end{enumerate}
檢查：搜尋 $0011$ 走 $1101 \to 0110$（第 $2$ 位 $0$，往左）$\to 0011 \to$ 自己，仍找得到。

\textbf{刪除 $1100$}。 搜尋 $1100$：標頭 $\to 1101$（第 $1$ 位 $1$，往右）$\to 1001$（第 $2$ 位 $1$，往右）
$\to 1100$（第 $4$ 位 $0$，往左）$\to$ 指回自己。 最後一個往回的指標是 $1100$ 自己的，屬於「$q = p$」的情形：
$1100$ 的另一個指標（右指標）指回 $1101$，把父節點 $1001$ 原本指向 $1100$ 的右指標改成這個往回的指標，
指向 $1101$，再釋放 $1100$。 $1101$ 本身的兩個指標沒有變。 檢查：搜尋 $1101$ 走 $1101 \to 1001$（第 $2$ 位 $1$，
往右）$\to$ 指回 $1101$，找得到。

\textbf{結果}。 $1101$ 的左指標指向 $0110$（插入時改的），右指標指向 $1001$（沒變）。 題目問 $1101$ 指向
哪些節點，選 (B)(D)。 (A) $0001$ 是標頭，指向 $1101$ 的是它；(C) $0011$ 原本是 $1101$ 的左子，插入後被 $0110$
隔開；(E) $1000$ 是 $1001$ 的左子。
\end{pastsoln}
